% 内容选择与规划后生成（Puduppully等，2019）的五个计算阶段及两项训练损失。
% 下方示例行使用正文的五条活动记录，为教学示例。
\begin{tikzpicture}[
  x=1mm,y=1mm,figure picture,
  stage/.style={figure node,minimum width=24mm,text width=22mm,minimum height=17mm,
    inner sep=2pt,font=\figurefont\scriptsize,align=flush center},
  ex/.style={font=\figurefont\scriptsize,text=FigureMuted,text width=26mm,align=flush center},
  loss/.style={figure node,draw=FigureLoss,fill=FigureLossFill,minimum height=6mm,
    inner sep=1.5pt,font=\figurefont\scriptsize},
  flow/.style={figure flow}
]
  \node[stage,draw=FigureInput,fill=FigureInputFill] (enc) at (0,0)
    {记录编码\\字段嵌入拼接后\\经前馈层得$\bm r_j$};
  \node[stage,draw=FigureModel] (gate) at (30,0)
    {内容选择门\\$\bm r^{\mathrm{cs}}_j=\bm g_j\odot\bm r_j$};
  \node[stage,draw=FigureModel] (plan) at (60,0)
    {指针规划\\$p(z_k=j\mid z_{<k},R)$};
  \node[stage,draw=FigureModel] (penc) at (90,0)
    {计划编码\\双向LSTM得$\bm e_k$};
  \node[stage,draw=FigureOutput,fill=FigureOutputFill] (dec) at (120,0)
    {文本解码\\注意力$\bm e_k$与复制};
  \foreach \a/\b in {enc/gate,gate/plan,plan/penc,penc/dec}{
    \draw[flow] (\a.east) -- (\b.west);
  }
  \node[ex,anchor=north] at (enc.south) {$r_1,\ldots,r_5$各为\\（实体，类型，值）};
  \node[ex,anchor=north] at (gate.south) {每条记录读取\\其他记录后得门值};
  \node[ex,anchor=north] at (plan.south) {$z=(r_1,r_2,r_3,r_5,$\\$\mathrm{EOS})$};
  \node[ex,anchor=north] at (penc.south) {只编码计划中的\\四条记录};
  \node[ex,anchor=north] at (dec.south) {“读书会于4月18日\\14:00在海棠楼201……”};
  \node[loss] (lp) at (60,14) {$-\log p_{\bm\phi}(z\mid R)$};
  \node[loss] (ly) at (120,14) {$-\log p_{\bm\theta}(y\mid R,z)$};
  \draw[figure link,densely dashed] (lp) -- (plan.north);
  \draw[figure link,densely dashed] (ly) -- (dec.north);
\end{tikzpicture}
